\section{ADIP Design}
\label{sec:design}

\subsection{Two scheduling layers}

Figure~\ref{fig:architecture} shows the paths under comparison. Direct mode
sends one OpenAI-compatible completion per client request. Pass-through adds
HTTP translation and telemetry but calls the backend once per request. Fixed
and adaptive modes place requests in compatibility-keyed queues and send one
backend prompt list per outer batch. In every mode, \vllm{} retains control of
its token-level continuous scheduler.

\begin{figure}[t]
  \centering
  \begin{tikzpicture}[
    node distance=10mm,
    box/.style={draw,rounded corners,align=center,minimum height=8mm,minimum width=24mm,font=\small},
    sched/.style={box,fill=orange!12},
    engine/.style={box,fill=blue!10},
    arrow/.style={-{Latex[length=2mm]},thick}
  ]
    \node[box] (client) {Open-loop\\client};
    \node[sched,right=of client] (gateway) {ADIP\\outer scheduler};
    \node[engine,right=of gateway] (vllm) {vLLM\\inner scheduler};
    \node[box,right=of vllm] (gpu) {RTX 3060\\GPU};
    \draw[arrow] (client) -- (gateway);
    \draw[arrow] (gateway) -- (vllm);
    \draw[arrow] (vllm) -- (gpu);
    \draw[arrow,bend right=28] (client.south) to node[below,font=\scriptsize]{direct mode} (vllm.south);
  \end{tikzpicture}
  \caption{The experimental boundary. Gateway batching controls when and in
  what HTTP grouping prompts become visible; \vllm{} independently controls
  token-level GPU scheduling.}
  \label{fig:architecture}
\end{figure}

\subsection{Compatibility and queue discipline}

A batch key consists of model identifier, output-token cap, and temperature.
Only equal keys share a backend request. The ingress queue is bounded; a full
queue rejects before creating unresolved work. One coordinator drains ingress
into per-key FIFO queues, selects the oldest compatible queue, obtains a pure
policy decision, and awaits one backend batch. This simple coordinator makes
the mechanism measurable, but it is intentionally serial: arrivals during a
backend call accumulate until that call returns. Section~\ref{sec:mechanism-model}
predicts that this occupancy-driven backlog can dominate the nominal
micro-wait.

\subsection{Fixed batch closure}

For the oldest compatible queue, the fixed policy dispatches when (i) the batch
reaches eight requests, (ii) the oldest request reaches the 1\,ms window, or
(iii) further waiting would consume estimated service slack before the earliest
soft deadline. Otherwise it waits only for the remaining window. Policy code
does not perform I/O or sleep; the coordinator owns timing and queue mutation.

\subsection{Adaptive closure}

The adaptive policy was frozen before final evaluation. It maintains
exponentially weighted moving averages of compatible arrival rate $\hat\lambda$
and backend time $\hat S$ with smoothing factor 0.2. For remaining window
$w_r$, it computes expected companions $\hat\lambda w_r$. If this value is
below 1.0, it dispatches immediately. Otherwise it waits for the minimum of
the remaining 20\,ms cap, estimated time to fill the batch,
\begin{equation}
  t_{\mathrm{fill}}=
  \frac{B_{\max}-B_{\mathrm{queue}}}{\hat\lambda},
\end{equation}
and positive deadline slack after reserving $\hat S$. Full queues, expired
windows, and exhausted slack dispatch immediately. Importantly, low-load
bypass does not split requests already accumulated during an in-flight backend
call; it closes the existing backlog without adding intentional wait.

\subsection{Instrumentation and failure semantics}

Every accepted request receives a UUID and optional experiment identifier.
Structured events record enqueue, policy reason, batch dispatch, backend
completion, and request completion without prompt text or credentials. The
response carries batch ID, batch size, queue time, backend time, total gateway
time, and deadline outcome. Backend failures resolve every future in the
affected batch. Cancellation and shutdown settle queued and in-flight futures
exactly once, preventing a successful test run from hiding leaked asynchronous
work.

The benchmark writes request JSONL as each request completes, while a separate
sampler records GPU utilization, VRAM, power, and temperature through
\texttt{nvidia-smi}. A manifest created before traffic binds the run to the
Git commit, model revision, mode, rate, seed, and frozen parameters. Atomic
summary writes occur only after traffic and validity checks complete.
